% arara: lualatex: { branch: developer, draft: true }
% arara: lualatex: { branch: developer, draft: true }
% arara: lualatex: { branch: developer }
% arara: clean: { extensions: [ aux, log, out] }
% arara: move: { files: [ enumext-09-luamml-mathml.html ], target: 'enumext-09-mathml.html' }
\DocumentMetadata
  {
    lang=en-US,
    pdfversion=2.0,
    pdfstandard=ua-2,
    tagging=on,
  }
\documentclass{article}
\usepackage[top=2cm,bottom=2cm,left=1cm,right=1cm]{geometry}
\usepackage{unicode-math,graphicx,hyperref,enumext}
\hypersetup
  {
    colorlinks = true,
    pdftitle   = {Test tagged PDF for enumext package},
  }
\begin{document}
\section{The environment \texttt{enumext}}

Default list.

\begin{enumext}
  \item One
  \item Two
  \item Three
\end{enumext}

Letters in the label.

\begin{enumext}[label=\alph*)]
  \item One
  \item Two
\end{enumext}

Four nested levels.

\begin{enumext}
  \item First level
    \begin{enumext}
      \item Second level
        \begin{enumext}
          \item Third level
            \begin{enumext}
              \item Fourth level
            \end{enumext}
        \end{enumext}
    \end{enumext}
  \item First level
\end{enumext}

Items with their own label and without label.

\begin{enumext}
  \item One
  \item[X] Two
  \item* Three
  \item Four
\end{enumext}

Start of the numbering.

\begin{enumext}[start=5]
  \item One
  \item Two
\end{enumext}

List without separation and text between items.

\begin{enumext}[nosep]
  \item One
    Text on the first item.
  \item Two
\end{enumext}

\section{The environment \texttt{enumext*}}

Horizontal list.

\begin{enumext*}
  \item One
  \item Two
  \item Three
  \item Four
\end{enumext*}

Letters in the label of the horizontal list.

\begin{enumext*}[label=\alph*)]
  \item One
  \item Two
  \item Three
\end{enumext*}

Horizontal list with three columns.

\begin{enumext*}[columns=3]
  \item One
  \item Two
  \item Three
  \item Four
  \item Five
\end{enumext*}

Joined items.

\begin{enumext*}[columns=3]
  \item One
  \item(2) Two and three
  \item Four
\end{enumext*}

Horizontal list inside a list.

\begin{enumext}
  \item First level
    \begin{enumext*}
      \item One
      \item Two
    \end{enumext*}
  \item First level
\end{enumext}

\section{Columns and minipage}

Two columns.

\begin{enumext}[columns=2]
  \item One
  \item Two
  \item Three
  \item Four
\end{enumext}

Minipage on the right.

\begin{enumext}[mini-env=0.4\linewidth]
  \item One
  \item Two
  \miniright
    Text in the minipage.
\end{enumext}

Columns and minipage.

\begin{enumext}[columns=2, mini-env=0.4\linewidth]
  \item One
  \item Two
  \item Three
  \item Four
  \miniright
    Text in the minipage.
\end{enumext}

Horizontal list and minipage.

\begin{enumext*}[columns=2, mini-env=0.4\linewidth,
                 mini-right={Text in the minipage.}]
  \item One
  \item Two
  \item Three
  \item Four
\end{enumext*}

\section{Fake \texttt{itemize} and \texttt{description}}

Fake itemize.

\begin{enumext}[nosep, label=\textbullet]
  \item First level item
    \begin{enumext}[nosep, label=\textendash]
      \item Second level item
        \begin{enumext}[nosep, label=\textasteriskcentered]
          \item Third level item
            \begin{enumext}[nosep, label=\textperiodcentered]
              \item Fourth level item
            \end{enumext}
        \end{enumext}
    \end{enumext}
  \item First level item
\end{enumext}

\newlength{\descitemwd}
\settowidth{\descitemwd}{\textbf{Something long}}

Fake description.

\begin{enumext}[label={}, labelwidth=\descitemwd, font=\bfseries,
                noitemsep, topsep=3pt]
  \item[Some] A short description.
  \item[] An entry without a label.
  \item[Something long] A longer description text may take more than one
    line or more than one paragraph.

    Another paragraph of the same item.
\end{enumext}

Fake description in a horizontal list.

\begin{enumext*}[label={}, labelwidth=\descitemwd, font=\bfseries,
                 noitemsep, topsep=3pt]
  \item[Some] A short description.
  \item[Something long] A longer description text.
\end{enumext*}

\section{Series, resume and reset}

Series at the first level.

\begin{enumext}[series=one]
  \item One
  \item Two
\end{enumext}

Text between the lists.

\begin{enumext}[resume=one]
  \item Three
  \item Four
\end{enumext}

Series at inner levels.

\begin{enumext}
  \item First level
    \begin{enumext}[series=foo]
      \item Second level
      \item Second level
    \end{enumext}
    Text on the first level.
    \begin{enumext}[resume=foo]
      \item Second level
    \end{enumext}
  \item First level
    \begin{enumext*}[series=baz]
      \item Second level starred
    \end{enumext*}
    Text on the first level.
    \begin{enumext*}[resume=baz]
      \item Second level starred
    \end{enumext*}
\end{enumext}

Resume without name and reset.

\begin{enumext}
  \item One
  \item Two
\end{enumext}

\begin{enumext}[resume*]
  \item Three
\end{enumext}

\begin{enumext}[reset]
  \item One again
\end{enumext}

\section{Answers with \texttt{anskey}}

Answers stored with \texttt{anskey}.

\begin{enumext}[save-ans=one, save-sep={ }, save-ref=true]
  \item Question A. \anskey{Answer A}
  \item Question B. \anskey{Answer B}
  \item Question C. \anskey[item-star]{Answer C}
\end{enumext}

Answer retrieved with \texttt{\textbackslash getkeyans}.

The answer to question \ref{one:2} is \getkeyans{one:2}.

Printed answers.

\printkeyans{one}

Printed answers in a horizontal list.

\printkeyans*{one}

Answers with \texttt{\textbackslash foreachkeyans}.

\foreachkeyans[sep={\par}]{one}

\section{Answers in environments}

Environment \texttt{anskey*}.

\begin{enumext}[save-ans=two]
  \item Question A.
    \begin{anskey*}
      Long answer with \verb+code+ and more text.
    \end{anskey*}
  \item Question B. \anskey{Answer B}
\end{enumext}

\printkeyans{two}

Environment \texttt{keyans}.

\begin{enumext}[save-ans=three]
  \item Question A.
    \begin{keyans}
      \item First
      \item* Second (correct)
      \item Third
    \end{keyans}
  \item Question B.
    \begin{keyans}
      \item*[some note] First (correct)
      \item Second
    \end{keyans}
\end{enumext}

\printkeyans{three}

Environment \texttt{keyans*}.

\begin{enumext}[save-ans=four]
  \item Question A.
    \begin{keyans*}
      \item First
      \item* Second (correct)
      \item Third
    \end{keyans*}
\end{enumext}

\printkeyans*{four}

Environment \texttt{keyans} with columns and minipage.

\begin{enumext}[save-ans=five]
  \item Question A.
    \begin{keyans}[mini-env=0.5\linewidth, columns=1]
      \item First
      \item* Second (correct)
      \item Third
      \miniright
        Text in the minipage.
    \end{keyans}
\end{enumext}

\section{Answers with images}

Environment \texttt{keyanspic}.

\begin{enumext}[save-ans=six, save-ref=true]
  \item Question with images.
    \begin{keyanspic}[layout-sty={3,2}]
      \anspic{\includegraphics[scale=0.15, alt={An uppercase letter A}]{example-image-a}}
      \anspic*{\includegraphics[scale=0.15, alt={An uppercase letter B}]{example-image-b}}
      \anspic{\includegraphics[scale=0.15, alt={An uppercase letter C}]{example-image-c}}
    \end{keyanspic}
  \item Question with images and label above.
    \begin{keyanspic}[label-pos=above]
      \anspic*[The letter A]{\includegraphics[scale=0.15, alt={An uppercase letter A}]{example-image-a}}
      \anspic{\includegraphics[scale=0.15, alt={An uppercase letter B}]{example-image-b}}
    \end{keyanspic}
\end{enumext}

\printkeyans*{six}

\section{Answers in horizontal lists}

Stored answers in \texttt{enumext*}.

\begin{enumext*}[save-ans=seven, save-sep={ }]
  \item Question A. \anskey{Answer A}
  \item Question B. \anskey{Answer B}
  \item Question C.
    \begin{anskey*}
      Answer C.
    \end{anskey*}
\end{enumext*}

\printkeyans{seven}

\section{Items without minipage}

\ExplSyntaxOn
\NewDocumentCommand \FILLER { m }
  { \prg_replicate:nn {#1} { This~sentence~repeats~to~fill~lines.~ } }
\ExplSyntaxOff

Full row items with \texttt{item-break}.

\begin{enumext*}[columns=2,item-break=true]
  \item One
  \item Two
  \item(2) Three, on the whole row.

    Second paragraph of the same item.
  \item Four
\end{enumext*}

\newpage

Without \texttt{item-break} a full row item stays in a \texttt{minipage}: it
cannot be split, so it moves as a whole to the next page and leaves a gap.

\vspace{0.8\textheight}

\begin{enumext*}[columns=2]
  \item One
  \item Two
  \item(2) \FILLER{40}
  \item Four
\end{enumext*}

\newpage

With \texttt{item-break=true} the same item starts on this page and continues
on the next one.

\vspace{0.8\textheight}

\begin{enumext*}[columns=2,item-break=true]
  \item One
  \item Two
  \item(2) \FILLER{40}
  \item Four
\end{enumext*}

Stored answers.

\begin{enumext}[save-ans=eight]
  \item Question A.
    \begin{keyans*}[columns=2,item-break=true]
      \item First
      \item Second
      \item(2)* Third (correct), on the whole row.
      \item Fourth
    \end{keyans*}
\end{enumext}

\printkeyans*{eight}

\section{Items without minipage and horizontal keys}

Each pair shows the same list with \texttt{item-break=true} and without it:
the lines of the full row item must start and end at the same places in both
lists, whatever the horizontal keys are.

\newcommand\KROWS
  {%
    \item One
    \item Two
    \item(2) Full row item with enough text to use two lines of the page: this
      sentence repeats. This sentence repeats to fill the line.

      Second paragraph of the full row item.
    \item Four
  }
\newcommand\KEYPAIR[2]
  {%
    \par\medskip\noindent\texttt{#2}, \texttt{item-break=true}:\par\nobreak
    \begin{#1}[columns=2,item-break=true,#2]
      \KROWS
    \end{#1}
    \par\medskip\noindent\texttt{#2}, without \texttt{item-break}:\par\nobreak
    \begin{#1}[columns=2,#2]
      \KROWS
    \end{#1}
  }

\KEYPAIR{enumext*}{list-offset=1cm}
\KEYPAIR{enumext*}{list-offset=-1cm}
\KEYPAIR{enumext*}{labelwidth=2cm,labelsep=5mm}
\KEYPAIR{enumext*}{itemindent=1cm}
\KEYPAIR{enumext*}{rightmargin=2cm}
\KEYPAIR{enumext*}{listparindent=1cm}
\KEYPAIR{enumext*}{align=right,labelwidth=1.5cm}
\KEYPAIR{enumext*}
  {list-offset=1cm,labelwidth=2cm,labelsep=5mm,rightmargin=1cm,itemindent=5mm}
% keyans* only exists inside enumext with save-ans, and needs an item*.
\newcommand\KEYPAIRK[1]
  {%
    \par\medskip\noindent\texttt{#1}, \texttt{item-break=true}:\par\nobreak
    \begin{enumext}[save-ans=keyx]
      \item Question.
        \begin{keyans*}[columns=2,item-break=true,#1]
          \item First
          \item Second
          \item(2)* Third (correct), on the whole row with enough text to use
            two lines of the page: this sentence repeats.

            Second paragraph of the full row item.
          \item Fourth
        \end{keyans*}
    \end{enumext}
    \par\medskip\noindent\texttt{#1}, without \texttt{item-break}:\par\nobreak
    \begin{enumext}[save-ans=keyy]
      \item Question.
        \begin{keyans*}[columns=2,#1]
          \item First
          \item Second
          \item(2)* Third (correct), on the whole row with enough text to use
            two lines of the page: this sentence repeats.

            Second paragraph of the full row item.
          \item Fourth
        \end{keyans*}
    \end{enumext}
  }
\KEYPAIRK
  {list-offset=1cm,labelwidth=2cm,labelsep=5mm,rightmargin=1cm,itemindent=5mm}

\par\medskip\noindent Alignment commands inside the full row item:

\begin{enumext*}[columns=2,item-break=true,labelwidth=2cm,labelsep=5mm]
  \item(2) {\centering Centered text, the indentation must be kept.\par}

    {\raggedright Text aligned to the left in the second paragraph.\par}
\end{enumext*}

\section{Defaults of the printed answers}

\texttt{\textbackslash setenumext} and \texttt{\textbackslash printkeyans}.

\setenumext[print,1]{columns=2, columns-sep=1cm, font=\bfseries}
\setenumext[print*]{columns=2, columns-sep=1cm, first=\bfseries}

\printkeyans{one}

\printkeyans*{one}
\end{document}
